\section*{第 \ref{ch:b2:anthropogenic-climate-change} 章——人为气候变化}

\begin{solution}{exo:b2:anthropogenic-climate-change:1}
(a) $3.75 - 0.60 + 0.22 + 0.45 + 0.11 = \qty{3.93}{W.m^{-2}}$，调整项合计
\qty{0.18}{W.m^{-2}}。(b) $3.93/3.71 = 1.059$：多 \qty{5.9}{\percent}。(c) 对流层
温度调整项 $-0.60$：对流层变暖后向太空发出更多红外线，带走了
部分多余能量，此时地表尚未作出响应。
\end{solution}

\begin{solution}{exo:b2:anthropogenic-climate-change:2}
(a) $4\pi(\num{6.371e6})^2 = \qty{5.10e14}{m^2}$。(b) 在 47 年内，$\num{434.9e21}/(47 \times
\num{3.156e7} \times \num{5.10e14}) = \qty{0.575}{W.m^{-2}}$；在 12 年内，$\num{152.4e21}/(12
\times \num{3.156e7} \times \num{5.10e14}) = \qty{0.789}{W.m^{-2}}$，与评估报告中的
0.57 和 0.79 相符。(c) $0.91 \times 434.9 = \qty{396}{ZJ}$。
\end{solution}

\begin{solution}{exo:b2:anthropogenic-climate-change:3}
(a) $-3.22 + 1.30 + 0.35 + 0.42 - 0.01 = \num{-1.16}$，即净值。(b) $3.93/1.16 =
\qty{3.39}{\celsius}$；$3.93/3.22 = \qty{1.22}{\celsius}$。(c) $3.39/1.22 = 2.78$：反馈
使普朗克增温几乎增至三倍。
\end{solution}

\begin{solution}{exo:b2:anthropogenic-climate-change:4}
(a) 1 拍克碳相当于 $44.009/12.011 = 3.664$~Gt 二氧化碳：$0.45 \times 3.664 =
\qty{1.65}{\celsius}$（每 \qty{1000}{PgC}），范围变为 0.989 至 2.31。(b) $2390 \times
12.011/44.009 = \qty{652}{PgC}$。(c) $0.45 \times 2.39 = \qty{1.08}{\celsius}$。
\end{solution}

\begin{solution}{exo:b2:anthropogenic-climate-change:5}
(a) $g = 159/398 = 0.400$；$T = (398/\num{5.670e-8})^{1/4} = \qty{289}{K}$。(b) $-0.6 \times 4
\times \num{5.670e-8} \times 288^3 = \qty{-3.25}{W.m^{-2}.K^{-1}}$；在 \qty{298}{K} 时，为 $-3.60$。
(c) 敏感度较低：响应随 $T^3$ 增长，因此温度更高的行星
在其他反馈不变时，更容易散去每增加一度所带来的额外能量。
\end{solution}

\begin{solution}{exo:b2:anthropogenic-climate-change:6}
(a) $1.85 - 0.50 = \qty{1.35}{W.m^{-2}.\celsius^{-1}}$。(b) $\sqrt{0.32^2 + 0.20^2} = 0.38$。
(c) 二者之和远比任一项更为确定：两者的误差呈强负相关。上对流层增温较多时，
其水汽含量也更多，因此较强的负\omterm{def:b2:anthropogenic-climate-change:feedbacks}{递减率反馈}
总是伴随着较强的正\omterm{def:b2:anthropogenic-climate-change:feedbacks}{水汽反馈}。
\end{solution}

\begin{solution}{exo:b2:anthropogenic-climate-change:7}
(a) $\lambda = (0.59 - 2.20)/1.03 = \qty{-1.56}{W.m^{-2}.\celsius^{-1}}$；$3.93/1.56 =
\qty{2.52}{\celsius}$。(b) $3.93/(1.56 - 0.5) = 3.93/1.06 = \qty{3.71}{\celsius}$。(c) $3.93/1.46 =
\qty{2.69}{\celsius}$。综合全部概率分布的评估结果为 3.5 和 \qty{2.7}{\celsius}。
\end{solution}

\begin{solution}{exo:b2:anthropogenic-climate-change:8}
(a) $0.617/3682 = \qty{1.68e-4}{\celsius^{-1}}$。(b) $0.113 \times 0.396 = \qty{0.0447}{m}$，即
\qty{44.7}{mm}，与预算中的 \qty{47.5}{mm} 接近。(c) $\num{2e-4} \times 0.2 \times 700 =
\qty{0.028}{m}$，即 \qty{28}{mm}。
\end{solution}

\begin{solution}{exo:b2:anthropogenic-climate-change:9}
(a) $\sum xy = 0.311$，$\sum x^2 = 0.300$：$\beta = 1.04$。(b) \qty{95}{\percent} 置信区间为
$1.04 \pm 1.96 \times 0.08$，即 0.88 至 1.20：区间不含 0，故响应被检测到；区间包含
1，故观测变化与模拟结果一致。(c) 区间 0.34 至 0.66 仍
不含 0（响应被检测到），但不含 1：模式会将响应高估约
两倍。
\end{solution}

\begin{solution}{exo:b2:anthropogenic-climate-change:10}
(a) $\sum g^2 = 1.20$，$\sum a^2 = 0.18$，$\sum ga = -0.46$，$\sum gy = 0.74$，$\sum ay = -0.28$。
正规方程为 $1.20\,\beta_g - 0.46\,\beta_a = 0.74$ 与 $-0.46\,\beta_g + 0.18\,\beta_a
= -0.28$；其行列式为 $1.20 \times 0.18 - 0.46^2 = 0.0044$，且 $\beta_g = (0.74 \times
0.18 - 0.46 \times 0.28)/0.0044 = 1.00$，$\beta_a = (-1.20 \times 0.28 + 0.46 \times 0.74)/0.0044 =
1.00$。(b) 温室气体 $+\qty{0.80}{\celsius}$，气溶胶 $-\qty{0.30}{\celsius}$。(c)
行列式很小，是因为 $g$ 与 $a$ 近似成正比（二者相关系数为 $-0.99$）：
$y$ 中的少量噪声可通过同时增大两个比例因子来吸收，因此二者各自
约束较差，尽管其合并效应已为人熟知。
\end{solution}

\begin{solution}{exo:b2:anthropogenic-climate-change:11}
(a) $(0.93 - 0.1)/0.45 \times 1000 = \qty{1844}{Gt}$；减去 210，并加上 $26 \times 0.93 = 24$：
\qty{1610}{Gt}。(b) $0.93 - x = 0.45 \times (1350 + 210 + 24)/1000 = 0.713$，故 $x =
\qty{0.22}{\celsius}$，接近评估报告中 0.1 至 \qty{0.2}{\celsius} 的上限。(c)
预算对非 CO$_2$ 增温 $x$ 的依赖，仅通过 $-x/\text{TCRE}$ 实现：每
\qty{0.1}{\celsius} 的代价为 $0.1/0.45 \times 1000 = \qty{222}{Gt}$，与限值无关。
\end{solution}

\begin{solution}{exo:b2:anthropogenic-climate-change:12}
(a) $3.93/1.81 = \qty{2.17}{\celsius}$，$3.93/0.51 = \qty{7.71}{\celsius}$，$3.93/1.16 =
\qty{3.39}{\celsius}$。(b) 低端长为 $3.39 - 2.17 = \qty{1.22}{\celsius}$，高端长为
$7.71 - 3.39 = \qty{4.32}{\celsius}$：函数 $1/|\lambda|$ 为凸函数，故
$\lambda$ 向零每移动等长一步，ECS 的增幅就越来越大。(c) 弱负的 $\lambda$ 给出极大的 ECS，
而 $\lambda$ 在零附近的微小变化难以排除；低 ECS 需要强负的
$\lambda$，\omterm{def:b2:anthropogenic-climate-change:feedbacks}{普朗克响应}对其设定了界限，而观测到的增温与之相矛盾。
\end{solution}

\begin{solution}{pb:b2:anthropogenic-climate-change:1}
\textbf{1.} $2.16 + 0.54 + 0.21 + 0.41 = \qty{3.32}{W.m^{-2}}$。
\textbf{2.} $-0.22 - 0.84 = \qty{-1.06}{W.m^{-2}}$。
\textbf{3.} $3.32 + 0.47 + 0.05 - 1.06 - 0.20 + 0.08 + 0.06 = \qty{2.72}{W.m^{-2}}$，即评估报告的
总量。
\textbf{4.} $3.93\log_2(410/278) = \qty{2.20}{W.m^{-2}}$，与 2.16 相差在 \qty{2}{\percent} 以内。
\textbf{5.} $3.93\log_2(427.35/278) = \qty{2.44}{W.m^{-2}}$。
\textbf{6.} 气溶胶–云项为 $-1.45$ 至 $-0.25$，范围宽 \qty{1.20}{W.m^{-2}}，而总量为 1.52：
它承担了总量的大部分不确定性，
因为其数值大且约束较差。
\textbf{7.} $0.79 \times \num{5.10e14} \times 12 \times \num{3.156e7} = \qty{1.53e23}{J}$，
\qty{153}{ZJ}。
\textbf{8.} $0.91 \times 153 = \qty{139}{ZJ}$。
\textbf{9.} $\lambda = (0.59 - 2.20)/1.03 = \qty{-1.56}{W.m^{-2}.\celsius^{-1}}$。
\textbf{10.} $3.93/1.56 = \qty{2.52}{\celsius}$。
\textbf{11.} $3.93 \times 1.03/2.20 = \qty{1.84}{\celsius}$。
\textbf{12.} $3.93/1.16 = \qty{3.39}{\celsius}$，高出 \qty{0.87}{\celsius}。
\textbf{13.} 迄今为止，增温在反馈最负的地区最强；随着
格局演变，反馈参数变得不那么负，因此长期敏感度
超过有效敏感度。
\textbf{14.} $0.45 \times 44.009/12.011 = \qty{1.65}{\celsius}$（每 \qty{1000}{PgC}）。
\textbf{15.} 0.989 至 \qty{2.31}{\celsius}（每 \qty{1000}{PgC}）。
\textbf{16.} $2390 \times 12.011/44.009 = \qty{652}{PgC}$。
\textbf{17.} $0.45 \times 2.39 = \qty{1.08}{\celsius}$。
\textbf{18.} 在舍入误差范围内，仅二氧化碳即可解释全部人为增温：其他
温室气体的增温与气溶胶的降温近乎相互抵消。
\textbf{19.} $1.5 - 1.07 = \qty{0.43}{\celsius}$。
\textbf{20.} $(0.43 - 0.1)/0.45 \times 1000 = \qty{733}{Gt}$ 二氧化碳。
\textbf{21.} $733 - 210 = \qty{523}{Gt}$。
\textbf{22.} $26 \times 0.43 = \qty{11.2}{Gt}$；$523 - 11.2 = \qty{512}{Gt}$。
\textbf{23.} 它减少 $0.1/0.45 \times 1000 = \qty{222}{Gt}$，降至 \qty{290}{Gt}：评估报告指出，
其他驱动因素的选择可使预算变动 \qty{220}{Gt}。
\textbf{24.} 由 TCRE 算得，自 2020 年初起，\qty{1.5}{\celsius} 对应的\omterm{def:b2:anthropogenic-climate-change:budget}{剩余碳预算}
为 \textbf{\qty[text-series-to-math=true]{512}{Gt}} 二氧化碳，与
评估报告在 \qty{50}{\percent} 概率下的 \qty{500}{Gt} 相差在 \qty{3}{\percent} 以内。
\end{solution}
